\chapter{Reproducibility Package}
\label{app:reproducibility}

This appendix describes the reproducibility infrastructure supporting the
experimental claims of this thesis.  All quantitative results reported in
Chapter~\ref{ch:experiments}---conformance rates, acceptance counts, prevention
detection rates, latency distributions, and statistical comparisons---are
derived from a versioned artifact store produced by a fully scripted,
deterministic pipeline~\citep{ammann2016testing,woodcock2009formalsurvey}.  The design separates \emph{data production} from
\emph{data presentation} so that each phase can be independently re-executed
and audited without rerunning costly language-model API calls.

The reproducibility package is organised into four logical \emph{artifact}
components that mirror the structure of a conventional empirical
software-engineering study:

\paragraph{Machine-readable schemas.}
The artifact releases \texttt{schemas/spec\_ir.schema.json} (canonical SpecIR)
and \texttt{schemas/semantic\_feedback.schema.json} (nine-field Semantic
Feedback IR).  The adapter registry (\texttt{src/adapters/\_\_init\_\_.py}) and
FSF lowerer (\texttt{src/ir/lowerers/fsf\_lowerer.py}) support round-trip
validation from SpecIR to legacy TaskSpec format.
\texttt{src/ir/schema\_validate.py} validates both schemas at test time via
\texttt{jsonschema} (see \texttt{tests/test\_schema\_validate.py}).

\paragraph{Raw execution logs.}
Append-only records of every pipeline invocation, one entry per
(task, mode, attempt) tuple.  Each record captures the prompt text
sent to the language model, the generated candidate source, formal
conformance scores, JSON-serialised Semantic Feedback IR payloads
(\texttt{semantic\_feedback} field), \texttt{attempt\_history} trajectories,
pattern hit lists with severity annotations, acceptance decisions, and
wall-clock timing.
Raw logs are the authoritative source of truth; all downstream
aggregates are derived from them.

\paragraph{Processed metrics.}
Deterministically aggregated summary tables at per-task, per-mode, and
per-run granularity.  Metrics include strict formal conformance,
binary acceptance, pattern violation counts partitioned by severity,
prevention detection rates, and latency percentiles.  Recomputing
processed metrics from identical raw logs yields bit-identical
outputs.

\paragraph{Figures and tables.}
Publication-quality visualisations and statistical summary tables
generated from processed metrics.  Each figure is linked to a manifest
entry specifying the data slice, filtering criteria, and research
question it addresses, enabling independent verification that the
plotted values correspond to the reported statistics.

\paragraph{Manuscript sources.}
The complete \LaTeX{} source tree for this thesis, including chapter
files, diagram sources, bibliography, and cross-references to
generated tables and figures.

This decomposition ensures that a sceptical reader can trace any reported
number from the manuscript back through processed metrics to the underlying
raw log entry without ambiguity.

% ---------------------------------------------------------------------------
\section{Running the Experiments}
\label{sec:repro:running}
% ---------------------------------------------------------------------------

Reproduction proceeds through four sequential pipeline stages.  Each stage consumes
the output of its predecessor and produces artefacts that are independently
archived.  Stages may be re-run in isolation provided their input artefacts
are available.

\paragraph{Stage 1: Benchmark construction.}
A deterministic task generator synthesises the hard benchmark suite
using a fixed pseudorandom seed.  The generator emits FSF
specifications with controlled difficulty parameters (input/output
arity, scenario count, guard overlap structure) and attaches
auto-generated reference implementations.  A satisfiability validator
discards tasks whose scenarios are individually unsatisfiable,
ensuring that every retained task admits at least one correct
implementation~\citep{demoura2008z3,barrett2018smt}.  The output is a frozen benchmark corpus used
identically across all experimental modes.

\paragraph{Stage 2: Experiment execution.}
The evaluation harness runs all compared pipeline modes over the
frozen benchmark.  For each (task, mode) pair, the harness invokes
the full synthesis--verification--repair loop (or the ablated variant
corresponding to the mode configuration), records every intermediate
result, and serialises the complete audit trace.  A separate prevention
evaluation protocol injects specification-layer and
implementation-layer mutants and measures detection rates~\citep{demillo1978mutation,li2023iet}.  This stage
is the only one that requires external language-model API access;
all subsequent stages operate on locally stored logs.

\paragraph{Stage 3: Data preparation.}
An aggregation pass reads raw execution logs and computes per-mode
summary statistics: mean and distribution of conformance scores,
acceptance rates, pattern hit frequencies, prevention detection and
false-accept rates, and paired comparison inputs for statistical
testing.  The aggregation is idempotent and seed-independent given
fixed raw logs.  Incomplete or interrupted runs are flagged rather
than imputed; corresponding table cells are marked \texttt{TBD} in
the manuscript.

\paragraph{Stage 4: Figure and table generation.}
Visualization and tabulation components consume processed metrics to
produce all figures and dynamically updated statistical tables
referenced in the manuscript.  Plot styling parameters (resolution,
colour palette, axis labels) are fixed to ensure visual reproducibility.
Statistical tests (Wilcoxon signed-rank, Mann--Whitney~U) are applied
at this stage with documented pairing structure and effect-size
computation.

As summarised in Figure~\ref{fig:repro-pipeline} (below), reproduction proceeds through
four sequential stages.

\begin{figure}[htbp]
  \centering
  \resizebox{\linewidth}{!}{%
  \begin{tikzpicture}[
    stage/.style={
      rectangle, rounded corners=3pt, draw=darkblue, thick,
      minimum width=2.6cm, minimum height=0.9cm, align=center,
      fill=blue!5, font=\small
    },
    arr/.style={-{Stealth[length=2.5mm]}, thick, darkblue}
  ]
    \node[stage] (s1) {Stage 1\\Benchmark construction};
    \node[stage, right=1.2cm of s1] (s2) {Stage 2\\Experiment execution};
    \node[stage, right=1.2cm of s2] (s3) {Stage 3\\Data preparation};
    \node[stage, right=1.2cm of s3] (s4) {Stage 4\\Figures \& tables};
    \draw[arr] (s1) -- (s2);
    \draw[arr] (s2) -- (s3);
    \draw[arr] (s3) -- (s4);
    \node[below=0.35cm of s1, font=\footnotesize\itshape] {Frozen benchmark};
    \node[below=0.35cm of s2, font=\footnotesize\itshape] {Raw logs};
    \node[below=0.35cm of s3, font=\footnotesize\itshape] {Processed metrics};
    \node[below=0.35cm of s4, font=\footnotesize\itshape] {Manuscript assets};
  \end{tikzpicture}%
  }
  \caption{Four-stage reproducibility pipeline.  Stages~1 and~2 produce
    primary data; Stages~3 and~4 derive all reported statistics and
    visualisations deterministically from stored artefacts.}
  \label{fig:repro-pipeline}
\end{figure}

\paragraph{Determinism and configuration control.}
Determinism is enforced at three levels to maximise reproducibility despite
the inherent stochasticity of language-model generation.

\paragraph{Task-level determinism.}
Benchmark construction uses a fixed pseudorandom seed, ensuring that every
replication operates on the identical task set with the same guard thresholds,
scenario orderings, and reference implementations.

\paragraph{Evaluation-level determinism.}
The formal conformance checker, SMT-based test-case generator, mutation
operators, and pattern guard are fully deterministic given a fixed candidate
and task.  Counterexample generation, branch-count analysis, and pattern
matching produce identical results across runs on the same hardware and
software stack.

\paragraph{Statistical-level robustness.}
Non-parametric statistical tests are employed for all mode comparisons because
they do not assume normality of conformance score distributions.  This choice
provides robustness to residual non-determinism in one-shot LLM generation
calls, where temperature settings above zero introduce sampling variability.
Repair-phase calls use zero temperature to approximate deterministic behaviour.

All experimental configuration parameters---mode definitions, synthesis budget
$K$, conformance threshold $\tau$, pattern severity thresholds, API model
identifier, token limits, and temperature settings---are recorded in a
machine-readable manifest bundled with the artifact store.  The manifest serves
as the single source of configuration truth for any replication attempt.

% ---------------------------------------------------------------------------
\section{Environment and Dependencies}
\label{sec:repro:env}
% ---------------------------------------------------------------------------

The experimental stack requires Python~3.10 or later, the Z3 SMT solver
(version~4.12 or compatible)~\citep{demoura2008z3}, and network access to an OpenAI-compatible
language-model API endpoint~\citep{openai2023gpt4}.  Core dependencies include abstract syntax tree
utilities (standard library), YAML parsing for pattern rule definitions, and
statistical libraries for non-parametric testing.  Optional plotting
dependencies are required only for Stage~4 figure generation.

Hardware requirements are modest: the formal checker and pattern guard are
CPU-bound and complete in milliseconds per candidate on a modern workstation.
The dominant cost is language-model API latency during Stage~2; the full
evaluation matrix (seven modes across 120 tasks) requires several hours of
API time but can be parallelised across independent worker processes without
affecting result validity.

% ---------------------------------------------------------------------------
\section{Expected Outputs}
\label{sec:repro:outputs}
% ---------------------------------------------------------------------------

A complete replication produces the following deliverables:

\begin{itemize}
  \item A frozen benchmark corpus with validated FSF specifications and
        reference implementations.
  \item Raw JSONL execution logs for every (task, mode) run, including
        prevention evaluation records.
  \item Processed metric summaries in structured JSON format, partitioned by
        mode and evaluation type (main experiment, spec-confusion prevention,
        implementation-screening prevention).
  \item Publication figures in vector and raster formats at 300\,dpi
        resolution.
  \item Updated statistical tables embedded in the manuscript source.
\end{itemize}

Incomplete runs are never silently extrapolated.  If a replication attempt
terminates before all tasks are processed, the aggregation stage reports the
actual completion count and marks pending entries as \texttt{TBD}.  This
conservative reporting policy prevents optimistic bias in acceptance rates and
statistical comparisons.

% ---------------------------------------------------------------------------
\section{Data Availability}
\label{sec:repro:data}
% ---------------------------------------------------------------------------

The artifact store, configuration manifest, and dependency specification are
retained alongside this thesis to enable independent verification of all
quantitative claims.  Raw execution logs are preserved in append-only form and
are not modified after initial write.  Processed metrics can be regenerated at
any time from the raw logs without re-invoking the language model.  Figure and
table assets are version-controlled together with the manuscript source so that
each reported statistic can be traced to its originating data slice.
